\documentclass[11pt]{article}
\usepackage{icmlko}
\begin{document}

\title{같은 손잡이가 다른 층에 있었다 --- 두 해 전 노트가 옳게 면제한 것을 내가 틀린 전제로 사전 등록했다}
\author{월드모델 실험실}
\date{노트 559 · 2026년 8월 5일}
\maketitle

\begin{abstract}
노트 $306$ 이 \textbf{위키}에서 \texttt{zero\_is\_data} 를 True$\to$False 로
정정했다 --- 마스크가 ``\emph{긁은 시점에 문서를 찾았나}''라 사후였다.
그 노트는 \textbf{검색을 명시적으로 면제}했다.

\begin{quote}
``\emph{검색 축은 이 문제가 없다. 빈 계열을 받은 것은 수집 시점의 진짜
측정이고, 마스크가 $0$ 인 것은 그 레코드의 \textbf{검색어를 아예 모른다}는
뜻이다.}''
\end{quote}

노트 $554$ 가 그 ``아예 모른다''가 \textbf{``라벨 하위 $73\%$ 다''}임을
보였다. 그래서 \textbf{노트 $306$ 자신의 검사를 검색에 돌렸다} --- 사전 수요
증거가 똑같이 없는 두 무리(긁혔는데 계열 $0$ · 아예 안 긁힘)의 라벨이
갈리는가.

\begin{center}
\begin{tabular}{lrrr}
\hline
 & $n_A$ & $n_B$ & $p$ \\
\hline
만화 학습 & $854$ & $4{,}747$ & $\mathbf{3.9 \times 10^{-293}}$ \\
펀딩 학습 & $295$ & $2{,}040$ & $2.3 \times 10^{-58}$ \\
애니 학습 & $275$ & $1{,}615$ & $3.4 \times 10^{-57}$ \\
게임 유보 & $29$ & $52$ & $0.83$ \\
\hline
\end{tabular}
\end{center}

\textbf{$7/10$ 칸에서 $p<0.01$ 로 갈린다.} 마스크가 수요가 아니라 수집을
나른다는 것은 확정됐다. \textbf{그런데 처방은 틀렸다.}
\end{abstract}

\begin{figure}[t]
\centering
\includegraphics[width=\textwidth]{figs/f1.pdf}
\caption{\textbf{가} --- 갈린다. \textbf{나} --- 손해가 나는 자리.
\textbf{다} --- 두 층.}
\end{figure}

\section{전역으로 고치니 값만 치른다}

\begin{center}
\begin{tabular}{lrrl}
\hline
자 (짝 씨앗 $12$) & 전 & 후 & 차 \\
\hline
\textbf{판} & $+0.4671$ & $+0.4637$ & $\mathbf{-0.0034}$ ($0/12$) \\
날짜 통제 판 & $+0.5041$ & $+0.5000$ & $-0.0040$ ($0/12$) \\
무리 안 판 & $+0.3868$ & $+0.3843$ & $-0.0025$ (문다) \\
KR 만화 & $+0.6390$ & $+0.6401$ & $+0.0011$ (문다) \\
비게임 앱 & $+0.2885$ & $+0.2858$ & $-0.0027$ (문다) \\
\hline
\end{tabular}
\end{center}

\textbf{손해가 어디서 나는지가 답이다.}

\begin{center}
\begin{tabular}{lrrr}
\hline
도메인 & 정정 전 유보 관측 & 차 & 무게 \\
\hline
\textbf{웹툰} & $\mathbf{99\%}$ & $\mathbf{-0.0216}$ & $19.3\%$ \\
\textbf{아이돌} & $\mathbf{100\%}$ & $\mathbf{-0.0552}$ & $1.5\%$ \\
게임 & $80\%$ & $-0.0137$ & $5.3\%$ \\
\textbf{도서} & $99\%$ & $\mathbf{+0.0459}$ & $4.8\%$ \\
\hline
만화 & $3\%$ & $+0.0007$ & 이미 \texttt{DOMDROP} \\
펀딩 & $0\%$ & $+0.0002$ & 이미 \texttt{BLOCK} \\
\hline
\end{tabular}
\end{center}

\section{배포 논증}

\textbf{배포에서는 새 작품을 반드시 긁는다.} 그러면
\texttt{zero\_is\_data=True} 의 마스크는 배포에서 \emph{항상} $1$ 이고,
\textbf{유보 관측이 $100\%$ 에 가까운 도메인은 이미 배포와 같다}
(웹툰 $99$ · 모바일 $100$ · 도서 $99$ · 아이돌 $100$). \textbf{거기엔 고칠
것이 없었다.}

\begin{quote}
그리고 진짜로 깨져 있던 곳 --- 만화($3\%$) · 펀딩($0\%$) · 애니 --- 은
\textbf{이미 \texttt{DOMDROP} 과 \texttt{BLOCK} 으로 처리돼 있다}. 전역
변경은 \textbf{안 깨진 데서 값을 치르고 깨진 데는 이미 고쳐져 있다.}
\end{quote}

\section{두 층은 다르다}

\begin{center}
\begin{tabular}{lll}
\hline
 & 위키(노트 $306$) & 검색(여기) \\
\hline
마스크의 뜻 & 지금 문서가 있나 & 우리가 긁었나 \\
\textbf{층} & \textbf{시점} 위반 --- 사후 & \textbf{선택} 인공물 --- 수집 순서 \\
처분 & 채택($0.4851 \to 0.4799$) & \textbf{기각} \\
\hline
\end{tabular}
\end{center}

\textbf{내 사전 등록은 둘을 같은 층으로 적었다} --- ``정정이므로 판이
내려가도 채택한다''. 그 문장의 ``정정이므로''가 \textbf{측정 전에 이미
결론을 담고 있었다.}

\begin{quote}
\textbf{노트 $548$ 과 같은 실수를 두 번째로 했다.} 그때는 ``이 값이
상한이다''라고 써 놓고 상한을 추정값처럼 쓰는 규칙을 달았고, 이번엔 ``이것은
정정이다''라고 써 놓고 정정임을 검증하지 않았다. \textbf{둘 다 사전 등록
문장 안에 검증되지 않은 전제를 박아 넣은 것이다.}
\end{quote}

\section*{반증된 예측 하나}

사전 등록 ①은 ``마스크$\leftrightarrow$라벨이 도메인마다 준다''였고
\textbf{반증됐다}(웹툰 $+0.059 \to +0.155$ · 게임 $+0.084 \to +0.174$ ·
도서 $-0.081 \to +0.145$). 만화만 $+0.2512 \to +0.0214$ 로 줄었다.

\begin{quote}
\textbf{방향을 잘못 적었다.} 위키는 ``지금 문서가 있나''(사후)$\to$``출시
전에 있었나''(사전)이라 상관이 \emph{줄고}, 검색은 ``우리가 긁었나''(수집)
$\to$``검색량이 문턱 위였나''(진짜 수요)라 \emph{올라야} 맞다. 노트 $306$
자신이 위키에서 남은 게임 $0.31$ 을 ``수선 뒤 진짜 사전 신호''라 적었다.
\end{quote}

\section*{조항}

\begin{quote}
\textbf{``이것은 정정이다''는 처분이지 관측이 아니다.} 사전 등록에 쓰려면
\emph{무엇이 정정인지}를 먼저 재야 한다 --- 여기서는 ``배포에서 그 마스크가
어떤 값인가''가 그 검사였고, 그것을 재니 고칠 데가 이미 고쳐져 있었다.
\end{quote}

\section*{남은 것}

\begin{itemize}
\item \textbf{도서만 False} --- 혼자 $+0.0459$ 로 크게 오른다(무게 $4.8\%$
$\to$ 판 기여 상한 $+0.0022$). 도메인별 \texttt{zero\_is\_data} 를 사전
등록해 재는 중(노트 $560$). \textbf{대조 팔}도 같이 돈다 --- 도서에서 검색을
\emph{통째로} 가리면 노트 $548$ 이 잰 대로 $-0.1654$ 여야 하고, 그러면
``값은 필요하고 관측-$0$ 만 해롭다''가 선다.
\item \textbf{세계애니 $82\%$ · 게임 $80\%$} --- 배포($100\%$)와 $20$\%p
떨어져 있는 둘. 노트 $554$ 의 자에서도 낮지 않다.
\item \textbf{아이돌 $-0.0552$} --- 관측이 $100\% \to 100\%$ 로 안 바뀌는데
왜 제일 크게 내려가나. 유보 $51$ 행이라 잡음일 수 있는데 안 봤다.
\end{itemize}

\end{document}
